\section{How the run proceeded}\label{sec:process}

\subsection{Timeline}
All times are UTC; the run started at 08:03 on 13 September 2026. Appendix~\ref{app:timeline} lists all 94 commits made during the run.

\paragraph{Hour 0--1 (triage).} The agent created the repository, extracted the text of the notebook, read the
new problems of the 21st issue itself (about 140 problems) and delegated the reading of issues 1--20 to five
sub-agents, each covering about 1500 lines of text and asked to return a compact list of problems ``settleable by
a finite computation or a short rigorous argument''. Two of the sub-agents exceeded their output limit and had to
be re-run with a request for a shorter report; part of one report (issues 19--20) was lost and this range was
re-triaged about seven hours later. In parallel the agent wrote the first search scripts (Problems 21.113, 21.29, 21.99,
21.52, 21.26, 21.115) and launched them in the background with log monitors.

\paragraph{Hour 1--6 (first results).} The first theoretical result, the decidability of Problem 21.32, was drafted at
08:52 and committed with its proof at 09:13 together with the solution of 16.60. The agent then worked through the
candidates from the triage: for each it wrote a GAP sweep over the SmallGroups library (or over simple groups,
primitive groups, character tables), started it, and moved on. By 13:00 about twenty sweeps were running. The
solutions of 2.78 (13:26), 18.46 (13:52) and 13.19 (14:21, generalised at 14:38) were found in this period, each
triggered by data from a sweep: for 18.46 the computation of minimal square-root overgroups showed that $D_8$
needs order 128; for 13.19 a search for the smallest configuration failed but suggested the class-2 construction
that the agent then wrote down and proved.

\paragraph{Hour 6--7 (compaction).} At about 14:40 the context of the main agent was compacted. The summary
preserved the state of the running computations and the list of open threads, and the agent continued from it. A
side effect surfaced at this point: the timestamps the agent had been writing into its progress log since about
13:00 were wrong by up to five hours (it had extrapolated the time instead of reading the clock). The error was
detected by comparing with the git history and corrected.

\paragraph{Hour 7--12 (the two big problems).} Problem 21.26 (Lisi--Sabatini) was verified up to order 2000 with a
new double-coset checker. Problem 16.14 (Berkovich) was attacked theoretically: the agent proved a structure
theorem for a minimal counterexample, realised that a counterexample
of order $2^{11}$ with $r=3$ would be a central extension of a group $H$ of order 256 with $d(H)=7$ satisfying a
cohomological condition, and excluded all ten such $H$; it then extended this to orders $2^{12}$ ($r=3$),
$2^9$ and $2^{10}$ ($r=2$). Problem 20.37 (Hooshmand)
was started at 15:50: the agent computed which factorisations of simple groups are obtainable from subgroup
chains, found that $\PSL(2,8)$ is the first group where this fails, proved a general theorem for the case of
two cosets (a perfect-matching argument in vertex-transitive graphs), and designed an exact-cover search for the
remaining cases. Problem 19.56 (Monakhov) was solved for minimal simple groups at 17:46, and 20.21
(Verret--Conder) received two impossibility lemmas.

\paragraph{Hour 12--40 (large sweeps).} The second half of the run was dominated by long computations: the
order-$2^{11}$ case of 16.14 required testing 7.1 million candidate quotients (about 6 hours on 14 cores, using the
$p$-covering group instead of cocycle computations, an idea found at hour 12); Problem 20.21 was searched
through all groups of order $\le 1000$ divisible by four (except two orders) and through the rank-$\le 4$ groups
of order 512; the 20.37 exact-cover searches ran in several phases with increasing effort. Meanwhile the agent
reviewed all reports line by line and corrected the wording of several claims.

\paragraph{Hour 40--48 and the end.} The agent planned the final consolidation for three hours before the
deadline but misread a 12-hour clock display and performed it at 09:44, 1\,h\,41\,min after the deadline. No
mathematical content was produced after the deadline: the final commit regenerated a table from data already on
disk, re-verified all 20.37 factorisations in GAP, stopped the remaining processes and added a closing entry to the
progress log.

\subsection{The working pattern}
Almost every problem was approached in the same way.
\begin{enumerate}
\item \emph{Formalise the statement as a finite check} (a GAP function \code{Check(G)} returning a certificate of
failure), test it on a few groups where the answer is known, and launch a sweep over a library in the background,
usually split into several processes by order or by id range.
\item \emph{Monitor} the log for hits and completions through the harness's monitor tool rather than by polling.
\item \emph{If a hit appears}, verify it with a second, independent computation (for 18.46 a lattice-based check
instead of \code{IsomorphicSubgroups}; for 13.19 a pc-presentation of the constructed group; for 20.37 a direct
multiplication of the sets in GAP after a Python search), then look for the general pattern behind it.
\item \emph{If no hit appears}, record the exact range verified, and think about whether the data suggests a
proof (as for 16.14, where the sweep showed that all inclusion-minimal Sylow intersections equal $O_p(G)$, or for
20.37, where the data showed which factor sizes need non-subgroup constructions).
\item \emph{Write the report as soon as a result is stable}, commit, and update the index.
\end{enumerate}
The theoretical work was interleaved with waiting for computations. The agent's own account of its reasoning
(preserved in the transcript) shows the typical pattern: a first idea, a concrete small case computed in GAP, a
revision of the idea, and finally a proof written directly into the report.

\subsection{Mistakes and how they were caught}
A record of the run would be incomplete without its errors; all of them were caught by the agent itself.
\begin{itemize}
\item \emph{Vacuous first version of a search} (20.122): the first sweep tested triples of Sylow subgroups only,
for which the question is trivial by Zenkov's theorem; the agent noticed this while writing the report and
replaced the sweep by one over all nilpotent subgroups (which is much more expensive, and was therefore only
completed for small orders).
\item \emph{Too strong a necessary condition} (20.21): the local condition was first formulated with an
automorphism of order exactly 3; the correct condition needs an automorphism of 3-power order, and the search was
repeated with the corrected condition.
\item \emph{A wrong side in a coset computation} (20.37): the first factorisation search used left cosets where
right cosets were required; the verification step (direct multiplication) rejected the ``solutions'' and the bug
was found immediately. This is the clearest example of the value of independent verification.
\item \emph{Clock drift in the progress log}, corrected from the git history (see above), and the late final
consolidation caused by a 12-hour clock display.
\item \emph{Operational accidents}: a process-killing command that matched its own shell (three times), duplicated
launches of the same sweep writing to the same log, GAP read-only global names, and a subagent report that exceeded
the output limit. None affected the results.
\item \emph{Over-optimistic literature claims}: the agent initially cited ``Konvisser's classification'' of
2-groups with three involutions as implying $d(G)\le 4$; it could not access the paper and later marked this as an
unverified literature dependency (it is now only needed for orders $\ge 2^{12}$ in Problem 16.14).
\end{itemize}

\subsection{Resource use}
The agent kept a running list of its processes and repeatedly stopped low-value sweeps to stay near the 20-core
budget (the load average was above 30 for short periods only). Memory was never an issue (peak use well below
20\,GB). The largest single computations were: the order-$2^{11}$ analysis for 16.14 (7.1 million groups, 6
hours on 14 cores); the order-512 rank-$\le 4$ search for 20.21 (420\,514 groups, about 8 hours on 8 cores); the
enumeration of all $1\,006\,700\,565$ connected graphs on 11 vertices for 20.71 (nauty, about 3 hours on 8
cores); and the 21.26 sweeps (about 800\,000 groups). The repository (without the session transcript and the
GAP installation) is about 530\,MB, mostly logs and intermediate data.
